%! @@TITLE@@ — Unit @@UNITINT@@, Lesson @@LESSONID@@ — @@DOCTITLE@@
\documentclass[10pt]{article}
\usepackage{@@PREFIX@@-article}
\usepackage{@@PREFIX@@-boxes}
\newcommand{\TallMath}[1]{$\displaystyle #1\rule[-1.4em]{0pt}{3.2em}$}

\begin{document}

\pageheader{Unit @@UNITINT@@, Lesson @@LESSONID@@}{@@DOCTITLE@@}
% NAMESTRIP: no \namedateperiod here. The student writes their name once, on the
% cover this component is stapled behind. See references/conventions.md.
@@NAMEROW@@
\vspace{0.15in}

% TODO: author the @@DOCTITLE@@. See references/components.md for the expected structure.
% If this is the HOMEWORK: it is this course's Check Your Understanding set — ~6 items in new
% contexts spanning the lesson's lettered skills, each ending in an interpretation or a
% justification, plus one spiral item and a closing \begin{spiralbox} preview. It IS scored.
% \boxguard before every breakable box — and in the _key too. Raise it (\boxguard[30]) when
% the box opens with an unbreakable TikZ figure or tabularx.
% Multi-step solutions go in a \begin{work} block authored BYTE-IDENTICALLY here and in the
% _key: the blank reserves its exact height, the key prints it. That is the work rule, and it
% is what keeps this component the same number of pages as its key.
\boxguard
\begin{notesbox}{@@DOCTITLE@@}
\writelines{4}
\end{notesbox}

\end{document}
